\bookchapter{บทที่ 16: ระบบบัญชีและบริการคลาวด์}

\chapterintro{ผมมีเพื่อนคนหนึ่ง ปีที่แล้วมือถือโดนขโมย สิ่งแรกที่เขานึกถึงไม่ใช่ “มือถือราคาเท่าไร” แต่เป็น “รูปในเครื่องจะเอากลับคืนมาได้ไหม” แล้วก็น่าเสียดาย เขาไม่ได้เปิดสำรองข้อมูลคลาวด์ มือถือเป็นแค่ของชิ้นหนึ่งข้างนอก แต่ข้อมูลในบัญชี หายแล้วหายเลย เรียกคืนไม่ได้อีกเลย บทนี้เราจะคุยกันว่า จะซื้อประกันที่แท้จริงให้ชีวิตดิจิทัลของคุณได้อย่างไร}

\booksection{หนึ่ง บัญชีมือถือของคุณ คุมอะไรมากแค่ไหนกันแน่?}

หลายคนคิดว่าบัญชีมีไว้โหลดแอปเท่านั้น ความจริงแล้ว ของที่ผูกอยู่หลังบัญชีมือถือหนึ่งใบ มีมากกว่าที่คุณคิดเยอะ:

\begin{itemize}
  \item \textbf{สมุดโทรศัพท์}: หมายเลขโทรศัพท์ที่คุณเก็บสะสมมาหลายปี หรือแม้แต่สิบกว่าปี
  \item \textbf{รูปภาพและวิดีโอ}: ตั้งแต่วันเกิดลูกจนถึงงานรับปริญญา จากเดตครั้งแรกถึงทริปฮันนีมูน
  \item \textbf{บันทึกย่อ}: ไอเดียที่จดสั้น ๆ ไดอารี รหัสผ่านต่าง ๆ (ใช่แล้ว หลายคนเก็บรหัสผ่านไว้ในบันทึกย่อ)
  \item \textbf{ประวัติซื้อ App}: แอปเสียเงินและการซื้อในแอปที่จ่ายไปหลายร้อย
  \item \textbf{ค้นหามือถือ}: พอมือถือหาย นี่คือเครื่องมือเดียวที่ช่วยคุณระบุตำแหน่ง ล็อกเครื่อง และล้างข้อมูลจากระยะไกลได้
\end{itemize}

ถ้าคุณใช้ iPhone ของทั้งหมดนี้จะผูกอยู่กับ Apple ID ของคุณ ถ้าใช้ Android ก็จะผูกกับบัญชี Huawei/Xiaomi/OPPO/vivo บวกกับบัญชี Google ของคุณ บัญชีเหล่านี้คือ “ทะเบียนบ้าน” ของชีวิตดิจิทัลคุณ ทะเบียนบ้านหาย เรื่องมากมายก็หมดสิทธิ์อ้าง

\booksection{สอง สามเรื่องที่ต้องทำ ทำเลยตอนนี้}

\begin{enumerate}
  \item \textbf{เปิดการยืนยันสองขั้น}: ไม่ว่าคุณจะใช้ iPhone หรือ Android เรื่องแรกคือเข้าไปที่การตั้งค่าบัญชี แล้วเปิด "การยืนยันสองขั้น" หรือ "การตรวจสอบสองขั้นตอน" ของตรงนี้เหมือนติดกลอนพิเศษให้ประตูบ้านคุณ — มีแค่รหัสผ่านไม่พอ ยังต้องให้มือถือของคุณรับรหัสยืนยันด้วย แม้ใครจะเดารหัสผ่านได้จากบนเน็ต แต่ถ้าไม่มีมือถือคุณ เขาก็เข้าไม่ได้
  \item \textbf{ผูกวิธีกู้คืนที่เชื่อถือได้เอาไว้อย่างหนึ่ง}: การยืนยันสองขั้นปลอดภัยมาก แต่ก็มีจุดเสี่ยง — ถ้ามือถือของคุณเองหายล่ะทำอย่างไร? ระบบจึงแนะนำให้ตั้ง "คีย์กู้คืน" หรือ "อีเมล/หมายเลขมือถือสำรอง" ตรงนี้คือทางเดียวที่จะเอาบัญชีคืนมาได้ในยามฉุกเฉิน ควรจดลงกระดาษมากที่สุด อย่าเก็บไว้ในมือถือ
  \item \textbf{ยืนยันว่าเปิด "ค้นหามือถือ" แล้ว}: ฟังก์ชันนี้ช่วยคุณระบุตำแหน่ง ล็อกเครื่อง เล่นเสียง และแม้แต่ล้างข้อมูลทั้งหมดจากระยะไกลได้ หลังมือถือหาย เป็นด่านสุดท้ายของการป้องกัน หลายคนกดข้ามขั้นตอนนี้ตอนซื้อเครื่องใหม่ พออยากใช้ทีหลังถึงบางอ้อว่าไม่ได้เปิด ก็สายเกินแก้
\end{enumerate}

\booksection{สาม iCloud: “อากาศหายใจ” ของผู้ใช้ Apple}

คนใช้ iPhone วันหนึ่งต้องเจอคำถามชวนปวดหัว: iCloud เต็มแล้วทำอย่างไรดี?

การสำรองข้อมูล

\begin{itemize}
  \item \textbf{Apple ให้ฟรี 5GB พูดตรง ๆ ว่า 5GB นี้แค่สำรองข้อมูลรอบเดียวยังไม่พอ ดังนั้นคนส่วนใหญ่สุดท้ายก็จะเจอป๊อปอัปน่ารำคาญนั่นคือ "พื้นที่ iCloud เต็มแล้ว" ตอนนั้นคุณมีสามทางเลือก:}ล้างพื้นที่
  \item \textbf{: เข้าการตั้งค่า ชื่อของคุณ iCloud จัดการพื้นที่จัดเก็บ ดูว่าอะไรกินพื้นที่คุณอยู่ โดยทั่วไปรูปภาพและการสำรองข้อมูลคือตัวหลัก สามารถเริ่มจากลบการสำรองข้อมูลเก่า ๆ ที่ไม่ต้องการออกได้}อัปเกรดพื้นที่
  \item \textbf{: 50GB เดือนละ 6 หยวน 200GB เดือนละ 21 หยวน ส่วนตัวคิดว่า 50GB สำหรับคนส่วนใหญ่ก็พอแล้ว ใช้หลัก ๆ สำรองข้อมูลกับเก็บรูปภาพ}ทางอ้อมที่ไม่ต้องเสียเงิน
\end{itemize}

: ถ้าไม่อยากจ่ายเงิน สามารถส่งออกรูปไปเก็บที่คอมพิวเตอร์เป็นประจำ แล้วปิดซิงก์รูป iCloud แต่เท่ากับว่าคุณจะเสียความสะดวกของการสำรองรูปอัตโนมัติไป

\begin{tcbinfo}{คำเตือนสำคัญอย่างหนึ่ง: คลังรูป iCloud ไม่ใช่ "เว็บไดรฟ์" — มันคือ "การซิงก์" คุณลบรูปหนึ่งใบบนมือถือ ฝั่ง iCloud ก็ลบตาม ถ้าสิ่งที่คุณต้องการคือ "สำรองขึ้นคลาวด์แล้วล้างเครื่อง" คลังรูป iCloud ไม่ใช่สิ่งที่คุณต้องการ สิ่งที่คุณต้องการคือ iCloud Backup หรือเว็บไดรฟ์ของผู้ให้บริการอื่น}
เพียงอัปเกรดเป็นแพ็กเกจ iCloud แบบชำระเงิน (50GB/200GB/2TB ฯลฯ) ก็จะกลายเป็น\textbf{iCloud+}ผู้ใช้ และปลดล็อกฟังก์ชันการปกป้องความเป็นส่วนตัว 3 อย่าง:

\begin{itemize}
  \item \textbf{Private Relay (รีเลย์ส่วนตัว)}: เวอร์ชัน "VPN" ของ Apple ทราฟฟิกเว็บผ่าน Safari จะถูกส่งผ่านเซิร์ฟเวอร์ของ Apple ก่อนไปยังเว็บไซต์ ไม่มีใคร (รวมถึง Apple) เห็นพร้อมกันทั้งสองอย่างว่าคุณคือใครและคุณกำลังเข้าชมอะไร เหมาะกับการปกป้องความเป็นส่วนตัวบน Wi-Fi สาธารณะ ตั้งค่าที่:\texttt{การตั้งค่า > ชื่อของคุณ > iCloud > รีเลย์ส่วนตัว}
  \item \textbf{Hide My Email (ซ่อนอีเมล)}: สร้างที่อยู่อีเมลสุ่มแบบใช้ครั้งเดียวเมื่อสมัครบนเว็บไซต์/App อีเมลจะถูกส่งต่อไปยังอีเมลจริงของคุณโดยอัตโนมัติ หากไม่ต้องการรับอีเมลจากเว็บไซต์นั้นอีก เพียงปิดที่อยู่สุ่มนี้ไปก็พอ อีเมลจริงของคุณจะไม่ถูกเปิดเผย เส้นทางการตั้งค่า:\texttt{การตั้งค่า > ชื่อของคุณ > iCloud > ซ่อนอีเมล}
  \item \textbf{โดเมนอีเมลแบบกำหนดเอง}: สามารถใช้โดเมนของคุณเอง (เช่น @นามสกุลของคุณ.com) ในการรับส่งอีเมล iCloud เหมาะสำหรับการสร้างแบรนด์ส่วนตัวหรือการจัดการอีเมลในครอบครัว
\end{itemize}

ข้อแนะนำ: Private Relay อาจใช้งานไม่ได้ในบางประเทศ/ภูมิภาค ผู้ใช้ในประเทศไทยอาจต้องใช้วิธีอื่นเสริม

\end{tcbinfo}

\begin{tcbwarn}{การปกป้องข้อมูลขั้นสูง iCloud (ADP) — การเข้ารหัสแบบ End-to-End อย่างแท้จริง}
ปลายปี 2022 Apple ได้เปิดตัว\textbf{การปกป้องข้อมูลขั้นสูง iCloud (Advanced Data Protection, ADP)}โดยขยายการเข้ารหัสแบบ End-to-End (E2EE) จากข้อมูลบางส่วนเดิม เช่น สุขภาพและรหัสผ่าน ไปยัง\textbf{ข้อมูล iCloud ส่วนใหญ่}—รวมถึงรูปภาพ, บันทึกย่อ, การสำรองข้อมูล iCloud, ข้อมูลอุปกรณ์ รวม 23 ประเภทข้อมูล

\begin{itemize}
  \item \textbf{E2EE คืออะไร}: คีย์เข้ารหัสจะถูกเก็บไว้ในอุปกรณ์ของคุณเท่านั้น แม้แต่เซิร์ฟเวอร์ของ Apple ก็ไม่มี แม้ว่า Apple จะถูกแฮ็กหรือหน่วยงานบังคับใช้กฎหมายขอเข้าถึงข้อมูล Apple ก็ไม่สามารถนำข้อมูลที่ถอดรหัสแล้วมาให้ได้
  \item \textbf{เส้นทางการเปิดใช้งาน}:\texttt{การตั้งค่า > ชื่อของคุณ > iCloud > การปกป้องข้อมูลขั้นสูง}แล้วทำตามคำแนะนำเพื่อตั้งค่าผู้ติดต่อกู้คืน/คีย์กู้คืน (\textbf{ต้องทำ}หากทำคีย์หาย ข้อมูลของคุณจะกู้คืนไม่ได้)
  \item \textbf{ข้อควรระวัง}: (1)หลังเปิดใช้งานแล้วต้องเก็บคีย์กู้คืนให้ดี ไม่เช่นนั้นหากบัญชีหายจะกู้คืนข้อมูลไม่ได้ (2)อัลบั้มรูปภาพที่แชร์, คลังภาพ iCloud ที่แชร์ และฟังก์ชันอื่นๆ อีกไม่กี่อย่างจะไม่เข้ารหัสแบบ E2EE (3)Apple ID ในภูมิภาคจีนแผ่นดินใหญ่ยังไม่รองรับ ADP ต้องสลับไปยังบัญชีภูมิภาคอื่น
\end{itemize}

คำแนะนำสั้นๆ: หากบัญชีของคุณมีข้อมูลส่วนตัวที่สำคัญ และภูมิภาคของบัญชีรองรับ ขอแนะนำอย่างยิ่งให้เปิดใช้งาน ADP นี่คือการปกป้องข้อมูลบนคลาวด์ที่ดีที่สุดที่คนทั่วไปจะได้รับ

\end{tcbwarn}

\booksection{สี่ ซิงก์กับการสำรองข้อมูล อย่าสับสน}

สองคำนี้ฟังดูคล้ายกัน แต่ผลลัพธ์ต่างกันคนละเรื่อง:

\begin{itemize}
  \item \textbf{ซิงก์}เป็นแบบสองทาง: คุณลบบนมือถือ คลาวด์ก็ลบ คุณลบบนคลาวด์ มือถือก็ลบตาม iCloud รูปภาพ สมุดโทรศัพท์ บันทึกย่อ ล้วนเป็นการซิงก์ ข้อดีคืออุปกรณ์หลายเครื่องให้สถานะตรงกัน ข้อเสียคือ "ลบแล้วหายจริง"
  \item \textbf{การสำรองข้อมูล (Backup)}เป็นแบบทางเดียว: มือถือถ่าย "สแนปช็ต" ของสถานะปัจจุบันเก็บขึ้นคลาวด์ หลังจากนั้นคุณจะลบอะไรบนมือถือ สำเนาในการสำรองข้อมูลบนคลาวด์ก็ยังอยู่ iCloud Backup การสำรองข้อมูลในเครื่องคอมพิวเตอร์ ล้วนเป็นการสำรองข้อมูล ข้อดีคือข้อมูลปลอดภัยขึ้น ข้อเสียคือกินพื้นที่ และไม่อัปเดตแบบเรียลไทม์
\end{itemize}

\booksection{ห้า เมทริกซ์บัญชีของผู้ใช้ Android}

ฝั่ง Android สถานการณ์ซับซ้อนกว่า Apple หน่อย เพราะข้อมูลของคุณอาจกระจายอยู่ในหลายบัญชี:

\begin{itemize}
  \item \textbf{บัญชีแบรนด์}(Huawei/Xiaomi/OPPO/vivo): คุมทั้งพื้นที่คลาวด์ ค้นหามือถือ ร้านแอป และกระเป๋าเงิน ทุกแบรนด์มีคลาวด์เป็นของตัวเอง ปกติเป็นโหมดฟรีพื้นฐาน + จ่ายเงินเพิ่มพื้นที่
  \item \textbf{บัญชี Google}: ถ้าคุณใช้ Gmail, Google Photos, Google สมุดโทรศัพท์ ของพวกนี้อยู่ในบัญชี Google ทั้งหมด Google ให้ฟรี 15GB จัดเต็มกว่า Apple แต่ก็ต้องแชร์กับ Gmail และ Drive
  \item \textbf{WeChat/QQ}: ทั้งสองเป็นอิสระจากระบบมือถือ ประวัติการแชทและรายการโปรดของ WeChat ไม่ได้อยู่ในบัญชีแบรนด์ ก็ไม่ได้อยู่ในบัญชี Google แต่เก็บอยู่บนเซิร์ฟเวอร์ของ WeChat เองเท่านั้น
\end{itemize}

ดังนั้นผู้ใช้ Android ตอนสำรองข้อมูล ต้องเช็กแยกแต่ละอย่าง: สำรองข้อมูลคลาวด์แบรนด์ สำรองข้อมูล Google และสำรองประวัติแชท WeChat ขาดอย่างใดอย่างหนึ่งไม่ได้

\booksection{หก การแชร์แบบครอบครัว: ให้คนในครอบครัวได้ประโยชน์ด้วย}

“การแชร์แบบครอบครัว” ของ Apple ให้คุณและคนในบ้านแชร์การซื้อ App เพลง หนัง พื้นที่ iCloud กันได้ ยังแชร์ตำแหน่งกันได้อีกด้วย เพิ่มได้มากสุด 6 คน วิธีตั้งค่าคือ: การตั้งค่า ชื่อของคุณ การแชร์แบบครอบครัว Huawei และ Xiaomi ก็มีฟังก์ชันคล้ายกัน หาได้ในการตั้งค่าบัญชีชื่อ “การแชร์แบบครอบครัว” หรือ “คลาวด์ครอบครัว”

สถานการณ์ที่ใช้ได้จริงมากอย่างหนึ่ง: ตั้งค่าการแชร์แบบครอบครัวไว้บนมือถือของคุณพ่อคุณแม่ แบบนี้คุณจะเห็นตำแหน่งของเขาได้ตลอด (โดยเขายินยอม) หากคุณที่เฒ่าหลงทางหรือติดต่อไม่ได้ ก็จะมีตำแหน่งให้ตามหา

\begin{tcbwarn}{ผลกรรมของบัญชีถูกขโมย ร้ายกว่าที่คุณคิด}
ถ้าแฮกเกอร์ได้รหัสผ่านบัญชีคุณไป เขาสามารถ: ล็อกมือถือคุณจากระยะไกลแล้วรีดไถเงินค่าไถ่ แอบดูรูปภาพและสมุดโทรศัพท์ของคุณ ใช้วิธีชำระเงินของคุณซื้อของ ที่ปวดหัวกว่านั้นคือ ถ้าบัญชีคุณผูกของไว้เยอะ เอาบัญชีคืนมาเพียงลำพังก็เป็นฝันร้ายเรื่องหนึ่งแล้ว ดังนั้น การยืนยันสองขั้น เปิดไปเลยตอนนี้

\end{tcbwarn}

\begin{tcbpractice}{ลงมือทำ: ทำการตรวจสุขภาพบัญชีของคุณให้ครบวงจรหนึ่งรอบ}
\begin{enumerate}
  \item : เปิดการตั้งค่ามือถือ เข้าไปที่หน้าบัญชี ยืนยันว่าเปิดการยืนยันสองขั้นแล้ว
  \item : เช็กว่า "หมายเลขโทรศัพท์ที่ไว้ใจ" และ "คีย์กู้คืน" ตั้งค่าไว้หรือยัง จดคีย์ลงกระดาษ
  \item : เข้าไปจัดการ iCloud หรือพื้นที่คลาวด์ ดูว่าข้อมูลอะไรกำลังซิงก์อยู่ ปิดตัวที่ไม่ต้องการ
  \item : ยืนยันว่าเปิด "ค้นหามือถือของฉัน" แล้ว
  \item : ถ้าคุณใช้ Android เช็กดูว่าการสำรองข้อมูลครบทั้งสามฝ่าย — คลาวด์แบรนด์ Google และ WeChat — หรือยัง
\end{enumerate}

\end{tcbpractice}
